\begin{lemma}
Assume the eventual comparisons
\(\noderef{ParameterHierarchyEventualComparisons}\), let
\[
K=\noderef{TwoBiteNaturalIndependenceScale}(1+\varepsilon,n),
\]
and suppose \(n>n_0\).  Conditioned on \(\mathcal D\), for every set \(I\)
with \(|I|=K\), the union of pairs
\[
\bigcup_{v\in L_I}\binom{X_v(I)}2
\]
has size at most \(\varepsilon_1 K^2\).
\end{lemma}
\begin{proof}
Fix \(I\) with \(|I|=K\).  Let
\[
\mathcal L=\{v:\noderef{TwoBiteLargeClass}(\omega,\varepsilon,I,v)\}
\]
be the large class, and for \(v\in\mathcal L\) write
\[
X_v=X_v(I)=\noderef{TwoBiteX}(\omega,I,v).
\]
The equality \(|I|=K\) is an equality of natural numbers with
\[
K=\noderef{TwoBiteNaturalIndependenceScale}(1+\varepsilon,n).
\]
Whenever \(K\) occurs in a real inequality below, it denotes the real cast
\((K:\mathbb R)=(|I|:\mathbb R)\), which is the scale in the Lean conclusion
\(\noderef{ClosedPairsBound}(\cdots,\varepsilon_1,(|I|:\mathbb R))\).
Unfolding \(\noderef{TwoBiteX}\) gives \(X_v=I\cap N_v\), so
\[
X_v\subseteq I. \tag{1}
\]
Unfolding \(\noderef{TwoBiteLargeClass}\), and using
\(\noderef{TwoBiteLargeCutoff}\) and \(\noderef{TwoBiteHugeCutoff}\), gives
\[
t_2:=n^{1/4+\varepsilon}<|X_v|
\le
t_1:=\frac{\sqrt{n\log n}}{\log\log n}. \tag{2}
\]

Unfolding \(\noderef{TwoBiteClosedPairsInClass}\), the set whose cardinality
appears in the Lean statement is
\[
\bigcup_{v\in\mathcal L}\binom{X_v}{2}.
\]
Thus, after unfolding \(\noderef{ClosedPairsBound}\), it is enough to prove
\[
\left|\bigcup_{v\in\mathcal L}\binom{X_v}{2}\right|
\le \varepsilon_1K^2. \tag{3}
\]

We next bound the total size of the sets \(X_v\).  The class \(\mathcal L\) is
a subfamily of
\(\noderef{TwoBiteLargeHugeVertices}(\omega,\varepsilon,I)\), since
large membership is the left disjunct in the definition of
\(\noderef{TwoBiteLargeHugeVertices}\).  Applying
\(\noderef{LargeHugeVertexCountBound}\) with the present
\(\noderef{FiberAndDegreeEvent}\), eventual comparisons, \(n>n_0\), and
\(|I|=K\), we get
\[
\ell:=|\mathcal L|<n^{1/4}. \tag{4}
\]
Explicitly, the child lemma bounds the real cardinality of
\(\noderef{TwoBiteLargeHugeVertices}(\omega,\varepsilon,I)\).  Since
\(\mathcal L\subseteq\noderef{TwoBiteLargeHugeVertices}(\omega,\varepsilon,I)\),
we first have
\[
(\ell:\mathbb R)\le
|\noderef{TwoBiteLargeHugeVertices}(\omega,\varepsilon,I)|,
\]
and then the displayed strict bound follows by transitivity.

For distinct \(u,v\in\mathcal L\), (1) gives
\[
X_u\cap X_v\subseteq N_u\cap N_v.
\]
We now apply \(\noderef{FiberAndDegreeEvent}\) to the right side of this
containment.  Set
\[
C_n:=100(\log n)^3.
\]
There are three cases.  If \(u=r_i\in V_R\) and \(v=r_j\in V_R\), then
distinctness of \(u\) and \(v\) gives \(i\ne j\), and the red-red lifted
intersection clause of \(\noderef{FiberAndDegreeEvent}\) gives
\[
|N_{r_i}\cap N_{r_j}|\le C_n.
\]
If \(u=b_i\in V_B\) and \(v=b_j\in V_B\), then \(i\ne j\), and the blue-blue
lifted intersection clause gives
\[
|N_{b_i}\cap N_{b_j}|\le C_n.
\]
Finally, if one of \(u,v\) is red and the other is blue, say
\(u=r_i\in V_R\) and \(v=b_j\in V_B\), then the mixed lifted intersection
clause gives
\[
|N_{r_i}\cap N_{b_j}|\le C_n;
\]
the reversed color order is the same statement with \(u\) and \(v\) swapped.
Combining the appropriate color-case bound with the containment above, and
casting cardinalities to real numbers, gives
\[
|X_u\cap X_v|\le C_n \tag{5}
\]
for every distinct \(u,v\in\mathcal L\).

Set
\[
S=\sum_{v\in\mathcal L}(|X_v|:\mathbb R),
\qquad
E=\sum_{\{u,v\}\in\binom{\mathcal L}{2}}
  (|X_u\cap X_v|:\mathbb R).
\]
Here \(\binom{\mathcal L}{2}\) denotes the unordered two-element subsets of
\(\mathcal L\); the summand is independent of the order in which the two
vertices are named because \(X_u\cap X_v=X_v\cap X_u\).
Enumerate \(\mathcal L=\{v_1,\ldots,v_\ell\}\) and write
\(Y_i=X_{v_i}\).  For this finite family,
\[
\left|\bigcup_{i=1}^{\ell}Y_i\right|
\ge
\sum_{i=1}^{\ell}|Y_i|
-
\sum_{1\le i<j\le \ell}|Y_i\cap Y_j|. \tag{6}
\]
Indeed, (6) follows by induction on the number of sets.  The step from
\(b\) sets to \(b+1\) sets uses
\[
|U\cup Y_{b+1}|=|U|+|Y_{b+1}|-|U\cap Y_{b+1}|,
\qquad U=\bigcup_{i=1}^{b}Y_i,
\]
and
\[
U\cap Y_{b+1}\subseteq
\bigcup_{i=1}^{b}(Y_i\cap Y_{b+1}),
\]
so
\[
|U\cap Y_{b+1}|
\le
\sum_{i=1}^{b}|Y_i\cap Y_{b+1}|.
\]
After casting all cardinalities in (6) to \(\mathbb R\), the first sum is
\(S\).  The second sum is \(E\), because the chosen enumeration gives a
bijection between index pairs \(1\le i<j\le\ell\) and unordered two-element
subsets \(\{v_i,v_j\}\) of \(\mathcal L\).  Thus
\[
\left(\left|\bigcup_{v\in\mathcal L}X_v\right|:\mathbb R\right)\ge S-E.
\]
The union is a subset of \(I\) by (1), hence
\[
K=(|I|:\mathbb R)\ge
\left(\left|\bigcup_{v\in\mathcal L}X_v\right|:\mathbb R\right)
\ge S-E. \tag{7}
\]
We now prove the stronger overlap estimate \(E\le S/2\), rather than bounding
\(E\) separately by \(K\).  Let \(\ell_{\mathbb N}=|\mathcal L|\) and let
\(\ell=(\ell_{\mathbb N}:\mathbb R)\).  The unordered pairs in
\(\binom{\mathcal L}{2}\) are counted by \(\ell_{\mathbb N}\choose 2\).  After
casting to \(\mathbb R\), \(\noderef{RealChooseTwo}\) gives
\[
(\ell_{\mathbb N}\choose 2:\mathbb R)
=\binom{\ell}{2}_{\mathbb R}
=\frac12\,\ell(\ell-1).
\]
Using the pairwise bound (5) term by term in the nonnegative finite sum
defining \(E\), we therefore have
\[
E\le \binom{\ell}{2}_{\mathbb R}C_n
=\frac12\,\ell(\ell-1)C_n.
\]
Because of (4), \(\ell<n^{1/4}\), hence
\[
\ell-1<n^{1/4}+\frac12.
\]
Also \(n_0<n\) implies \(1\le n\), so \(\log n\ge0\), and therefore
\[
C_n=100(\log n)^3\ge0.
\]
The eventual-comparisons package contains
the monotonicity-size comparison
\[
100\left(n^{1/4}+\frac12\right)(\log n)^3
\le \noderef{TwoBiteLargeCutoff}(\varepsilon,n)=t_2.
\]
Thus
\[
(\ell-1)C_n
\le \left(n^{1/4}+\frac12\right)100(\log n)^3
\le t_2. \tag{8}
\]
Since each \(v\in\mathcal L\) satisfies \(t_2<|X_v|\) by (2), (8) gives the
pointwise inequality
\[
(\ell-1)C_n<|X_v|
\qquad\text{for every }v\in\mathcal L.
\]
Summing this pointwise inequality over the finite set \(\mathcal L\) and using
the previous bound on \(E\), we get
\[
2E\le \ell(\ell-1)C_n
=\sum_{v\in\mathcal L}(\ell-1)C_n
\le \sum_{v\in\mathcal L}|X_v|=S.
\]
The middle equality is just the real cast of the natural identity saying that
there are \(\ell_{\mathbb N}\) terms in the constant sum over \(\mathcal L\).
Therefore
\[
E\le \frac12S. \tag{9}
\]
Combining (7) and (9) yields
\[
K\ge S-E\ge S-\frac12S=\frac12S,
\]
and hence
\[
S\le 2K. \tag{10}
\]

Now bound the closed pairs.  Since
\(\noderef{TwoBiteClosedPairsInClass}\) is the finite union, over
\(v\in\mathcal L\), of the pair sets \(\binom{X_v}{2}\),
\[
\left|\bigcup_{v\in\mathcal L}\binom{X_v}{2}\right|
\le
\sum_{v\in\mathcal L}\binom{|X_v|}{2}_{\mathbb R}.
\]
This is only the finite union bound.  No disjointness of the pair sets is
assumed: if the same unordered base pair lies in \(\binom{X_u}{2}\) and
\(\binom{X_v}{2}\), it is counted once on the left and once in each relevant
summand on the right.  For each fixed \(v\), the real cardinality of
\(\binom{X_v}{2}\) is \(\binom{|X_v|}{2}_{\mathbb R}\), by
\(\noderef{RealChooseTwo}\), so summing these individual cardinalities bounds
the real cardinality of the union.
For each \(v\), the upper half of (2) gives
\[
\binom{|X_v|}{2}_{\mathbb R}
\le \frac12\,t_1|X_v|.
\]
This is a real-valued calculation: if \(x=|X_v|\), then
\[
\binom{x}{2}_{\mathbb R}=\frac{x(x-1)}2\le\frac{x^2}{2}
\le\frac{t_1x}{2},
\]
because \(x\ge0\) and \(x\le t_1\) by (2).  Therefore, using (10),
\[
\left|\bigcup_{v\in\mathcal L}\binom{X_v}{2}\right|
\le
\frac12t_1S
\le t_1K. \tag{11}
\]
Finally, \(\noderef{ParameterHierarchyEventualComparisons}\), applied to this
\(n>n_0\), contains the explicit rounded inequality
\[
t_1K\le\varepsilon_1K^2.
\]
Combining this with (11) proves (3), which is precisely
\(\noderef{ClosedPairsBound}\) for the cardinality of
\(\noderef{TwoBiteClosedPairsInClass}(\omega,I,
\noderef{TwoBiteLargeClass}(\omega,\varepsilon,I))\) at scale \(K=|I|\).
\end{proof}
